\section{Thiết kế Model Deployment}

\subsection{Lưu trữ và phân phối model}

Model sau tối ưu (Chương~\ref{chap:toi-uu}) được lưu trên một kho model kiểu
Hugging Face Hub, có thể là huggingface.co, mirror hoặc máy chủ tự dựng. Mỗi
model được tham chiếu bằng \texttt{ModelRef(repo\_id, filename, revision)}.
Khi phân giải, \texttt{HuggingFaceModelRepository} chuyển \texttt{revision}
(nhánh hoặc tag) thành \emph{commit SHA} cụ thể và sinh URL tải đã ghim theo SHA
đó (\texttt{ModelLocation}). Vì URL trỏ đến nội dung bất biến của một commit, hai
node triển khai cùng một deployment chắc chắn nhận cùng một tệp model, kể cả khi
nhánh \texttt{main} được cập nhật giữa chừng.

Phía Edge, \texttt{HttpArtifactFetcher} tải model qua HTTPS theo từng khối 1~MB
vào tệp tạm \texttt{.part}, sau đó đổi tên nguyên tử thành tệp đích. Tệp được lưu
cache theo hàm băm SHA-256 của URL (bỏ phần query), nên triển khai lại cùng model
không tải lại. URI \texttt{file://} hoặc đường dẫn cục bộ được dùng trực tiếp,
thuận tiện cho kiểm thử trên thiết bị.

\subsection{Luồng triển khai}

Luồng triển khai đầy đủ được mô tả trong Hình~\ref{fig:model-deployment-design}.

\begin{figure}[H]
\centering
\resizebox{\textwidth}{!}{%
\begin{tikzpicture}[font=\small,
  msg/.style={font=\scriptsize, above, fill=white, inner sep=1pt},
  note/.style={font=\scriptsize, align=left, anchor=west, draw=gray, fill=yellow!8, rounded corners=2pt, inner sep=2pt}]
  \foreach \x/\n/\t in {0/u/{Dashboard hoặc Assistant}, 3.4/c/{Cloud API}, 6.8/m/{MQTT}, 10.2/a/{Edge Agent}, 13.6/r/{Docker, AI Runtime}}
    { \node[draw, rounded corners, fill=blue!5, align=center, text width=2.6cm, minimum height=1.1cm] (\n) at (\x,0) {\t};
      \draw[dashed, gray] (\x,-0.6) -- (\x,-11.2); }
  \draw[flowarrow] (0,-1.3) -- node[msg]{POST /deployments} (3.4,-1.3);
  \node[note] at (3.5,-2.0) {lưu Runtime (CREATING),\\Deployment (IN\_PROGRESS)};
  \draw[flowarrow] (3.4,-2.9) -- node[msg]{command/deploy} (6.8,-2.9);
  \draw[flowarrow, dashed] (3.4,-3.5) -- node[msg]{202 Accepted} (0,-3.5);
  \draw[flowarrow] (6.8,-4.1) -- node[msg]{command/deploy} (10.2,-4.1);
  \node[note] at (10.3,-4.9) {chống trùng, lưu desired state,\\tải model (có cache)};
  \draw[flowarrow] (10.2,-5.9) -- node[msg]{pull image, run} (13.6,-5.9);
  \draw[flowarrow, dashed] (13.6,-6.6) -- node[msg]{container id} (10.2,-6.6);
  \draw[flowarrow] (10.2,-7.3) -- node[msg]{event/deployment} (6.8,-7.3);
  \draw[flowarrow] (6.8,-7.9) -- node[msg]{DeploymentEvent} (3.4,-7.9);
  \node[note] at (3.5,-8.7) {apply\_result: SUCCEEDED,\\Runtime RUNNING};
  \draw[flowarrow] (10.2,-9.6) -- node[msg]{state (retained)} (6.8,-9.6);
  \draw[flowarrow] (6.8,-10.2) -- node[msg]{NodeState} (3.4,-10.2);
  \draw[flowarrow, dashed] (3.4,-10.8) -- node[msg]{WebSocket} (0,-10.8);
\end{tikzpicture}}
\caption{Trình tự triển khai model từ Dashboard đến AI Runtime}
\label{fig:model-deployment-design}
\end{figure}

\subsection{Cập nhật và khôi phục}

Cập nhật model hoặc cấu hình là một lệnh \texttt{update} với spec mới trên cùng
\texttt{runtime\_id}. Agent chỉ tạo lại container khi spec thực sự thay đổi. Nếu
khởi chạy với model mới thất bại, Agent báo FAILED kèm thông điệp lỗi (ví dụ
không tải được model, image không tồn tại); người vận hành hoặc Assistant có thể
gửi lại một \texttt{update} trỏ về model trước đó để khôi phục. Vì mỗi lệnh
update mang \texttt{command\_id} mới, kết quả của lần thử lỗi trước không ghi đè
kết quả khôi phục. Việc tự động canary, chia đợt và rollback theo ngưỡng telemetry
được xác định là hướng phát triển (Chương~\ref{chap:ket-luan}).
